\section*{第 \ref{ch:b3:complex-kinetics} 章——复杂动力学：链、酶与振荡}

\begin{solution}{exo:b3:complex-kinetics:1}
$\ce{Cl2 -> 2 Cl}$：链引发；$\ce{Cl + H2 -> HCl + H}$ 和 $\ce{H + Cl2 -> HCl + Cl}$：
链增长；$\ce{2 Cl -> Cl2}$：链终止。\omterm{def:b3:complex-kinetics:chain}{链载体}为 Cl 和 H。总反应（两个
链增长步骤之和）：$\ce{H2 + Cl2 -> 2 HCl}$。
\end{solution}

\begin{solution}{exo:b3:complex-kinetics:2}
$[\mathrm M]_{1/2} = k_2/k_{-1} = \qty{1.0e-3}{mol/L} = \qty{1.0}{mol/m^3}$; $p = cRT = 1.0 \times 8.314 \times 750 =
\qty{6.2}{kPa}$.
\end{solution}

\begin{solution}{exo:b3:complex-kinetics:3}
$v/V_{\max} = [\mathrm S]/(K_M + [\mathrm S])$：$\frac12$、$\frac34$、$\frac{9}{10}$。90\,\% 的 $V_{\max}$ 需要 $[\mathrm S] = 9K_M$。
\end{solution}

\begin{solution}{exo:b3:complex-kinetics:4}
$k_{\mathrm{cat}}/K_M = 100/\num{0.80e-3} = \qty{1.25e5}{L.mol^{-1}.s^{-1}}$：比扩散
极限低约 60\,000 倍，接近典型值。
\end{solution}

\begin{solution}{exo:b3:complex-kinetics:5}
\ce{CH3CO} 的稳态条件：$k_2[\ce{CH3}][\ce{CH3CHO}] = k_3[\ce{CH3CO}]$。\ce{CH3} 的稳态条件：
$k_1[\ce{CH3CHO}] - k_2[\ce{CH3}][\ce{CH3CHO}] + k_3[\ce{CH3CO}] - 2k_4[\ce{CH3}]^2 = 0$。两式相加得 $k_1[\ce{CH3CHO}] =
2k_4[\ce{CH3}]^2$，故 $[\ce{CH3}] = (k_1/2k_4)^{1/2}[\ce{CH3CHO}]^{1/2}$，且 $\dd[\ce{CH4}]/\dd t = k_2[\ce{CH3}][\ce{CH3CHO}] =
k_2(k_1/2k_4)^{1/2}[\ce{CH3CHO}]^{3/2}$.
\end{solution}

\begin{solution}{exo:b3:complex-kinetics:6}
转化率为一半时，$[\ce{H2}] = [\ce{Br2}] = \frac12$，$[\ce{HBr}] = 1$（相对单位）：
$v/v_0 = \frac12 \times (\frac12)^{1/2}/(1 + 0.10 \times 1/0.5) = 0.354/1.2 = 0.29$。无抑制时
则为 0.354：抑制使速率再除以 1.2。
\end{solution}

\begin{solution}{exo:b3:complex-kinetics:7}
爆炸发生的条件为 $1000p > 6000/p + 10p^2$，即 $10p^3 - 1000p^2 + 6000 < 0$，其正
根为 2.48 和 \qty{99.9}{kPa}：混合物在约 2.5 与
\qty{100}{kPa} 之间发生爆炸（模型的第一、第二\omterm{def:b3:complex-kinetics:branching}{爆炸极限}）。
\end{solution}

\begin{solution}{exo:b3:complex-kinetics:8}
$1/v = 1/V_{\max} + (K_M/V_{\max})/[\mathrm S]$：$5.0 = 1/V_{\max} + 2.0K_M/V_{\max}$ 和 $2.5 = 1/V_{\max} + 0.5K_M/V_{\max}$。
两式相减得 $K_M/V_{\max} = \qty{1.667}{s.mM.\micro M^{-1}}$，于是 $1/V_{\max} = 1.667$：$V_{\max} = \qty{0.60}{\micro M.s^{-1}}$，
$K_M = \qty{1.0}{mM}$.
\end{solution}

\begin{solution}{exo:b3:complex-kinetics:9}
在 $1/v$ 坐标轴上截距相同，为\omterm{def:b3:complex-kinetics:inhibition}{竞争性抑制}。$K_M^{\mathrm{app}} = 1/0.417 = \qty{2.40}{mM} = 3K_M$，故
$1 + 50/K_i = 3$，从而 $K_i = \qty{25}{\micro M}$。
\end{solution}

\begin{solution}{exo:b3:complex-kinetics:10}
$N = \qty{1.01}{mol/L}$：$t_{1/2} = \ln(1.01/0.010 - 1)/(0.50 \times 1.01) = \ln100/0.505 = \qty{9.1}{s}$。
速率 $kx(N - x)$ 在 $x = N/2$ 时最大，
即同一时刻，也是 S 形曲线的拐点。
\end{solution}

\begin{solution}{exo:b3:complex-kinetics:11}
迹为 $B - 2$，行列式为 1。$B = 1.5$：$\lambda = -0.25 \pm 0.968\iu$，为稳定焦点（阻尼
振荡趋向稳态）。$B = 2.5$：$\lambda = 0.25 \pm 0.968\iu$，为不稳定焦点：
振荡逐渐增长，直至达到\omterm{def:b3:complex-kinetics:oscillating}{极限环}。
\end{solution}

\begin{solution}{exo:b3:complex-kinetics:12}
$1/\tau = \num{1.0e8} \times \num{5.40e-5} + \num{1.0e3} = \qty{6.4e3}{s^{-1}}$，$\tau = \qty{156}{\micro s}$。在几种
总浓度下测量 $\tau$：$1/\tau$ 对 $[\mathrm A]_e + [\mathrm B]_e$ 作图为直线，斜率为 $k_1$，
截距为 $k_{-1}$。
\end{solution}

\begin{solution}{pb:b3:complex-kinetics:1}
\textbf{1.} 起始时 $[\mathrm S]$ 已知，而可能抑制反应或发生逆向反应的
产物尚不存在。
\textbf{2.} 低 $[\mathrm S]$ 的数据点，即 $1/[\mathrm S]$ 与 $1/v$ 较大的点：它们决定斜率，而
取倒数会放大它们的误差。
\textbf{3.} 它按原始测量值对各速率进行加权；莱恩威弗–伯克直线给出的
参数有偏倚。
\textbf{4.} $K_M$：$0.773$ 比 0.801 低 1.2 个标准误；$V_{\max}$：$0.491$ 比
0.502 低 2.2 个标准误。两者均偏低，与偏倚的预期相符。
\textbf{5.} $k_{\mathrm{cat}} = 0.502/0.0050 = \qty{100}{s^{-1}}$.
\textbf{6.} $100/\num{0.801e-3} = \qty{1.25e5}{L.mol^{-1}.s^{-1}}$.
\textbf{7.} 约低于扩散极限 60\,000 倍；是典型的酶。
\textbf{8.} \omterm{def:b3:complex-kinetics:inhibition}{竞争性抑制}：$V_{\max}$ 不变，表观 $K_M$ 随 $[\mathrm I]$ 增大。
\textbf{9.} $K_M^{\mathrm{app}} = K_M(1 + [\mathrm I]/K_i) = K_M + (K_M/K_i)[\mathrm I]$.
\textbf{10.} $K_i = 0.799/0.0321 = \qty{24.9}{\micro M}$.
\textbf{11.} $24.9 \pm 4.30 \times 0.9 = 24.9 \pm \qty{3.9}{\micro M}$.
\textbf{12.} 酶–抑制剂复合物的解离常数，即
占据一半游离酶的抑制剂浓度。
\textbf{13.} $v_i/v_0 = (K_M + [\mathrm S])/(K_M(1 + [\mathrm I]/K_i) + [\mathrm S])$.
\textbf{14.} $2/3 = 0.67$.
\textbf{15.} $2/12 = 0.17$.
\textbf{16.} $11/21 = 0.52$：底物浓度高时，底物
竞争胜过抑制剂。
\textbf{17.} 当 $[\mathrm S] = K_M$ 时，$v_i/v_0 = 2/(2 + [\mathrm I]/K_i) = 0.10$，由此得 $[\mathrm I]/K_i = 18$。
\textbf{18.} $18 \times 24.9 = \qty{448}{\micro M}$.
\textbf{19.} 第一步很快：每个酶分子迅速释放一个 $\mathrm P_1$
并被捕获为酰基酶；此后 $\mathrm P_1$ 仅
以慢的第二步释放酶的速率生成。
\textbf{20.} $1.28/2.0 = 0.64 = (k_2/(k_2 + k_3))^2$，故 $k_2/(k_2 + k_3) = 0.80$。
\textbf{21.} $200/2.0 = \qty{100}{s^{-1}}$，与第一部分相同。
\textbf{22.} $k_2 = 0.80 \times 625 = \qty{500}{s^{-1}}$, $k_3 = \qty{125}{s^{-1}}$.
\textbf{23.} $500 \times 125/625 = \qty{100}{s^{-1}}$.
\textbf{24.} \textbf{$[\mathrm I]_{90\,\%} = 18K_i \approx \qty{0.45}{mM}$，在 $[\mathrm S] = K_M$ 时。}
\end{solution}
